\subsection{通过标量扩张得到的模的对偶}

设 A、B 为两个环， \(\rho: \mathbf{A} \to \mathbf{B}\) 一个环同态，E 是一个左 A-模，且 \(\mathbf{E}^*\) 其对偶。我们将定义一个典范的B-线性映射。

\[
\nu _ {E}: (E ^ {*}) _ {(B)} \rightarrow (E _ {(B)}) ^ {*}.\tag{19}
\]

(19) 的左边可写成 \(\mathrm{Hom}_{\mathbf{A}}(\mathbf{E},\mathbf{A})\otimes_{\mathbf{A}}\rho_{*}(\mathbf{B}_{s})\)，其中在 \(\mathrm{Hom}_{\mathbf{A}}(\mathbf{E},\mathbf{A})\) 中，\(\mathbf{A}\) 被视为 \((\mathbf{A},\mathbf{A})\)-双模。则存在一个典范 \(\mathbf{Z}\)-同态（\(\S 4\)，第2号，公式(7)）

\[
\nu: \operatorname{Hom} _ {A} (E, A) \otimes_ {A} \rho_ {*} (B _ {s}) \rightarrow \operatorname{Hom} _ {A} (E, A \otimes_ {A} \rho_ {*} (B _ {s})) = \operatorname{Hom} _ {A} (E, \rho_ {*} (B _ {s}))
\]

利用§3第4号命题4中的典范同构给出的等同关系。另一方面，(19)式的右端可以写成 \(\mathrm{Hom}_{\mathbf{B}}(\rho_{*}(\mathbf{B}_{d}) \otimes_{\mathbf{A}} \mathbf{E}, \mathbf{B}_{s})\) ；由于B是(B, A)-双模，存在一个典范的Z-同构（§4第1号命题1）

\[
\beta : \operatorname{Hom} _ {\mathbf {B}} (\rho_ {*} (\mathbf {B} _ {d}) \otimes_ {\mathbf {A}} \mathbf {E}, \mathbf {B} _ {s}) \rightarrow \operatorname{Hom} _ {\mathbf {A}} (\mathbf {E}, \operatorname{Hom} _ {\mathbf {B}} (\mathbf {B} _ {s}, \mathbf {B} _ {s}))
\]

和 \(\mathrm{Hom}_{\mathbb{B}}(\mathbf{B}_s,\mathbf{B}_s)\) 作为 A-模，典范地等同于 \(\rho_{*}(\mathbf{B}_{s})\) （参见第 1 号，命题 1 的证明）。考虑到这些等同，我们得到同态 \(\nu_{\mathbb{E}}\) ；容易验证该同态由等式

\[
\langle \xi \otimes x, \nu _ {E} (x ^ {*} \otimes \eta) \rangle = \xi_ {\rho} (\langle x, x ^ {*} \rangle) \eta ,\tag{20}
\]

（对 \(x \in \mathbf{E}\)、\(x^{*} \in \mathbf{E}^{*}\)、\(\xi\)、\(\eta\) 属于 \(\mathbf{B}\)）所刻画，这立即表明 \(\nu_{\mathbf{E}}\) 是B-线性的。此外，对每个A-线性映射 \(u: \mathbf{E} \to \mathbf{F}\)，图表

\[
\begin{array}{c c c} (\mathrm{F} ^ {*}) _ {(\mathrm{B})} & \xrightarrow {\nu_ {\mathrm{F}}} & (\mathrm{F} _ {(\mathrm{B})}) ^ {*} \\ \Big \downarrow^ {(t _ {u}) _ {(\mathrm{B})}} & & \Big \downarrow^ {t (u _ {(\mathrm{B})})} \\ (\mathrm{E} ^ {*}) _ {(\mathrm{B})} & \xrightarrow {\nu_ {\mathrm{E}}} & (\mathrm{E} _ {(\mathrm{B})}) ^ {*} \end{array}
\]

是交换的。

命题8. 若A-模E与 \(\rho_{*}(B_{s})\) 之一是投射且有限生成的，则同态 \(\nu_{\mathbb{E}}\) 是双射。

这由上述内容及第4节第2号命题2得出。

特别地，假设E是有限生成的自由A-模，令 \((e_{i})_{1\leqslant i\leqslant n}\) 为E的一组基，\((e_{i}^{*})\) 为其对偶基；则典范同构（19）把基 \((e_{i}^{*}\otimes1)\)（属于 \((\mathrm{E}^{*})_{(\mathrm{B})}\)）映为 \((1\otimes e_{i})\)（\(\mathrm{E}_{(\mathrm{B})}\) 的基）的对偶基。

